\documentclass[10pt,a4paper]{amsart}
\usepackage{amsmath,amssymb,amsthm,graphicx}
\usepackage{verbatim}
\newenvironment{reference}{\comment}{\endcomment}
%\newenvironment{reference}{}{}
\newenvironment{slogan}{\comment}{\endcomment}
\newenvironment{history}{\comment}{\endcomment}

% For commutative diagrams we use Xy-pic
\usepackage[all]{xy}

% We use 2cell for 2-commutative diagrams.
\xyoption{2cell}
% Theorem environments.
%
\theoremstyle{plain}
\newtheorem{theorem}[subsection]{Theorem}
\newtheorem{proposition}[subsection]{Proposition}
\newtheorem{lemma}[subsection]{Lemma}

\theoremstyle{definition}
\newtheorem{definition}[subsection]{Definition}
\newtheorem{example}[subsection]{Example}
\newtheorem{exercise}[subsection]{Exercise}
\newtheorem{situation}[subsection]{Situation}

\theoremstyle{remark}
\newtheorem{remark}[subsection]{Remark}
\newtheorem{remarks}[subsection]{Remarks}

\numberwithin{equation}{subsection}

% Macros
%
\def\lim{\mathop{\mathrm{lim}}\nolimits}
\def\colim{\mathop{\mathrm{colim}}\nolimits}
\def\Spec{\mathop{\mathrm{Spec}}}
\def\Hom{\mathop{\mathrm{Hom}}\nolimits}
\def\Ext{\mathop{\mathrm{Ext}}\nolimits}
\def\SheafHom{\mathop{\mathcal{H}\!\mathit{om}}\nolimits}
\def\SheafExt{\mathop{\mathcal{E}\!\mathit{xt}}\nolimits}
\def\Sch{\mathit{Sch}}
\def\Mor{\mathop{\mathrm{Mor}}\nolimits}
\def\Ob{\mathop{\mathrm{Ob}}\nolimits}
\def\Sh{\mathop{\mathit{Sh}}\nolimits}
\def\NL{\mathop{N\!L}\nolimits}
\def\CH{\mathop{\mathrm{CH}}\nolimits}
\def\proetale{{pro\text{-}\acute{e}tale}}
\def\etale{{\acute{e}tale}}
\def\QCoh{\mathit{QCoh}}
\def\Ker{\mathop{\mathrm{Ker}}}
\def\Im{\mathop{\mathrm{Im}}}
\def\Coker{\mathop{\mathrm{Coker}}}
\def\Coim{\mathop{\mathrm{Coim}}}

% Boxtimes
%
\DeclareMathSymbol{\boxtimes}{\mathbin}{AMSa}{"02}

%
% Macros for moduli stacks/spaces
%
\def\QCohstack{\mathcal{QC}\!\mathit{oh}}
\def\Cohstack{\mathcal{C}\!\mathit{oh}}
\def\Spacesstack{\mathcal{S}\!\mathit{paces}}
\def\Quotfunctor{\mathrm{Quot}}
\def\Hilbfunctor{\mathrm{Hilb}}
\def\Curvesstack{\mathcal{C}\!\mathit{urves}}
\def\Polarizedstack{\mathcal{P}\!\mathit{olarized}}
\def\Complexesstack{\mathcal{C}\!\mathit{omplexes}}
% \Pic is the operator that assigns to X its picard group, usage \Pic(X)
% \Picardstack_{X/B} denotes the Picard stack of X over B
% \Picardfunctor_{X/B} denotes the Picard functor of X over B
\def\Pic{\mathop{\mathrm{Pic}}\nolimits}
\def\Picardstack{\mathcal{P}\!\mathit{ic}}
\def\Picardfunctor{\mathrm{Pic}}
\def\Deformationcategory{\mathcal{D}\!\mathit{ef}}
\usepackage{cleveref}
\crefname{originallemma}{original lemma}{original lemmas}
\Crefname{originallemma}{Original lemma}{Original lemmas}
\makeatletter
\newlabel{lemma-characterize-smooth-over-field}{{00TT}{source}{}{sourcepointer-1}{}}
\newlabel{lemma-characterize-smooth-over-field@cref}{{[originallemma][0][]00TT}{[1][0][]source}{}{}{}}
\makeatother
\crefname{originallemma}{original lemma}{original lemmas}
\Crefname{originallemma}{Original lemma}{Original lemmas}
\makeatletter
\newlabel{lemma-etale}{{00U2}{source}{}{sourcepointer-2}{}}
\newlabel{lemma-etale@cref}{{[originallemma][0][]00U2}{[1][0][]source}{}{}{}}
\makeatother
\crefname{originallemma}{original lemma}{original lemmas}
\Crefname{originallemma}{Original lemma}{Original lemmas}
\makeatletter
\newlabel{lemma-smooth-at-generic-point}{{07ND}{source}{}{sourcepointer-3}{}}
\newlabel{lemma-smooth-at-generic-point@cref}{{[originallemma][0][]07ND}{[1][0][]source}{}{}{}}
\makeatother
\crefname{originallemma}{original lemma}{original lemmas}
\Crefname{originallemma}{Original lemma}{Original lemmas}
\makeatletter
\newlabel{proposition-dimension-zero-ring}{{00KJ}{source}{}{sourcepointer-4}{}}
\newlabel{proposition-dimension-zero-ring@cref}{{[originallemma][0][]00KJ}{[1][0][]source}{}{}{}}
\makeatother
\crefname{originallemma}{original lemma}{original lemmas}
\Crefname{originallemma}{Original lemma}{Original lemmas}
\makeatletter
\newlabel{theorem-nullstellensatz}{{00FV}{source}{}{sourcepointer-5}{}}
\newlabel{theorem-nullstellensatz@cref}{{[originallemma][0][]00FV}{[1][0][]source}{}{}{}}
\makeatother
\title{An algebra over a field is etale exactly when it is a finite product of finite separable extensions}
\author{The Stacks Project Authors}
\date{stacks-project:algebra.tex:local-snapshot-20261012; GNU FDL 1.2 or later}
\begin{document}
\maketitle
\noindent\textit{Editorial note.} All delivered mathematical text is the author's own, copied from the pinned Stacks Project chapter file algebra.tex; no statement or proof step was rewritten, reordered or removed.\par\medskip
\noindent\textit{Editorial note.} Preamble provenance, as required for this source: the frozen snapshot stages the chapter's own preamble file beside it (preamble.tex, 217 lines, 8116 bytes, the same shared Stacks-wide preamble the published chow and etale-cohomology cards use); the first line of the pinned chapter file is \textbackslash{}input\{preamble\}, so no preamble text lives inside the chapter itself. The card takes that file's lines 8-18 (the verbatim package and the three \textbackslash{}comment-based environments reference/slogan/history at lines 9-12, section 13.2, together with the Xy-pic setup) and its lines 154 to the runtime end of the file (the native theorem declarations and mathematical macros). Not one character was changed, no package was added, no macro was defined, and the excerpt serves only as preamble, never as a span coordinate.\par\medskip
\noindent\textit{Editorial note.} The same-chapter pointers are located at their \textbackslash{}label lines inside the pinned chapter file; the cross-chapter pointers are located in the project's own tag list staged in the snapshot (sources/stacks-chow/tags.txt), at the verbatim lines carrying their canonical tags.\par\medskip
\noindent\textit{Editorial note.} The counted proof refers to same-chapter statements by label; they are recorded as navigation stubs pointed at their \textbackslash{}label lines, and the arguments stay inside the counted span.\par\medskip
\noindent\textit{Editorial note.} The counted spans contain no \textbackslash{}cite, so the card carries no bibliography.\par\medskip
\noindent\textit{Allowed standard prerequisites (editorial).}\begin{itemize}
\item Etaleness is stable under finite products. Used at the start of the counted proof (lines 39618-39621). Reduces the statement to the individual factors of a finite product.
\item Smoothness at the generic point forces separability. Invoked at the end of the counted proof (lines 39646-39648). Shows the field extension is separable once the algebra is smooth.
\end{itemize}\medskip
\ifdefined\hypertarget\hypertarget{sourcepointer-1}{}\fi
\noindent\textit{Source reference.} Original lemma 00TT: sources/stacks-algebra/pinned-original-source.tex, source lines 38826--38826. Navigation pointer for \textbackslash{}ref\{lemma-characterize-smooth-over-field\} used inside the counted span (line 39636); the argument that consumes it is retained in the span.\par\medskip
\ifdefined\hypertarget\hypertarget{sourcepointer-2}{}\fi
\noindent\textit{Source reference.} Original lemma 00U2: sources/stacks-algebra/pinned-original-source.tex, source lines 39507--39507. Navigation pointer for \textbackslash{}ref\{lemma-etale\} used inside the counted span (line 39621); the argument that consumes it is retained in the span.\par\medskip
\ifdefined\hypertarget\hypertarget{sourcepointer-3}{}\fi
\noindent\textit{Source reference.} Original lemma 07ND: sources/stacks-algebra/pinned-original-source.tex, source lines 39108--39108. Navigation pointer for \textbackslash{}ref\{lemma-smooth-at-generic-point\} used inside the counted span (line 39646); the argument that consumes it is retained in the span.\par\medskip
\ifdefined\hypertarget\hypertarget{sourcepointer-4}{}\fi
\noindent\textit{Source reference.} Original lemma 00KJ: sources/stacks-algebra/pinned-original-source.tex, source lines 14632--14632. Navigation pointer for \textbackslash{}ref\{proposition-dimension-zero-ring\} used inside the counted span (line 39639); the argument that consumes it is retained in the span.\par\medskip
\ifdefined\hypertarget\hypertarget{sourcepointer-5}{}\fi
\noindent\textit{Source reference.} Original lemma 00FV: sources/stacks-algebra/pinned-original-source.tex, source lines 6741--6741. Navigation pointer for \textbackslash{}ref\{theorem-nullstellensatz\} used inside the counted span (line 39644); the argument that consumes it is retained in the span.\par\medskip
\begin{lemma}
\label{lemma-etale-over-field}
Let $k$ be a field. A ring map $k \to S$ is \'etale if and only if $S$
is isomorphic as a $k$-algebra to a finite product of finite separable
extensions of $k$.
\end{lemma}

\begin{proof}
We are going to use without further mention: if
$S = S_1 \times \ldots \times S_n$ is a finite product
of $k$-algebras, then $S$ is \'etale over $k$ if and only
if each $S_i$ is \'etale over $k$. See Lemma \ref{lemma-etale} part (11).

\medskip\noindent
If $k'/k$ is a finite separable field extension then we can
write $k' = k(\alpha) \cong k[x]/(f)$. Here $f$ is the minimal
polynomial of the element $\alpha$. Since $k'$ is separable over $k$
we have $\gcd(f, f') = 1$. This
implies that $\text{d} : k'\cdot f \to k' \cdot \text{d}x$
is an isomorphism. Hence $k \to k'$ is \'etale. Thus if $S$
is a finite product of finite separable extensions of $k$, then
$S$ is \'etale over $k$.

\medskip\noindent
Conversely, suppose that $k \to S$ is \'etale. Then $S$ is smooth
over $k$ and $\Omega_{S/k} = 0$. By
Lemma \ref{lemma-characterize-smooth-over-field}
we see that $\dim_\mathfrak m \Spec(S) = 0$ for every
maximal ideal $\mathfrak m$ of $S$. Thus $\dim(S) = 0$. By
Proposition \ref{proposition-dimension-zero-ring}
we find that $S$ is a finite product of Artinian local rings.
By the already used
Lemma \ref{lemma-characterize-smooth-over-field}
these local rings are fields. Hence we may assume $S = k'$ is a field.
By the Hilbert Nullstellensatz (Theorem \ref{theorem-nullstellensatz})
we see that the extension $k'/k$ is finite. The smoothness
of $k \to k'$ implies by Lemma \ref{lemma-smooth-at-generic-point}
that $k'/k$ is a separable extension and
the proof is complete.
\end{proof}

\end{document}
